\section{Set Theory and B Mathematics}
\label{sec:set-theory-b-mathematics}
\label{sec:b-core-language}

The core language used in this thesis is a typed expression and predicate language.
Following the usual B presentation, predicates are expressions of boolean type.
This keeps the syntax uniform: logical connectives, membership tests, equalities, and quantified formulas are all terms.

\subsection{Types}

The formalized type language is deliberately small:
\[
  \tau, \sigma ::= \intB
  \mid \boolB
  \mid \setB \tau
  \mid \tau \prodB \sigma.
\]
The types \(\intB\) and \(\boolB\) classify integer and boolean values.
The type \(\setB \tau\) classifies sets whose elements have type \(\tau\).
The product type \(\tau \prodB \sigma\) classifies ordered pairs.
There is no primitive function type in this B core.
As in B, a partial function from a set \(S : \setB \tau\) to a set \(T : \setB \sigma\) is a set of pairs satisfying a functionality condition.
Consequently, the expression \(S \pfunB T\) has type
\[
  \setB(\setB(\tau \prodB \sigma)).
\]
It denotes the set of all partial functions whose graph is included in \(S \times T\).

This design is important for the SMT encoding.
An SMT backend tends to offer native function types, but B functions are mathematical objects: they are sets of pairs, have domains and ranges, can be members of other sets, and can be compared as values.
Later chapters explain why translating such objects directly to SMT functions is convenient but requires care.

\subsection{Terms and predicates}

The grammar below summarizes the core terms.
The metavariables \(x,y,z\) range over variables, \(n\) over integers, \(b\) over booleans, and \(v_1,\ldots,v_k\) over a nonempty list of bound variables.
\[
\begin{array}{rcl}
t,u &::=& x
  \mid n
  \mid b
  \mid t \mapletB u
  \mid t + u
  \mid t - u
  \mid t \cdot u
  \mid t \leB u
  \mid t \andB u
  \mid \notB t
  \mid t \eqB u
  \\[0.2em]
&& \mid \mathbb{Z}
  \mid \mathbb{B}
  \mid t \inB u
  \mid \{v_1,\ldots,v_k \inB D \mid P\}
  \mid \powB t
  \mid t \prodB u
  \mid t \cupB u
  \mid t \capB u
  \mid |t|
  \\[0.2em]
&& \mid t(u)
  \mid \lambda v_1,\ldots,v_k \inB D.\,t
  \mid t \pfunB u
  \mid \min(t)
  \mid \max(t)
  \mid \forall v_1,\ldots,v_k \inB D.\,P.
\end{array}
\]
The constants \(\mathbb{Z}\) and \(\mathbb{B}\) denote the sets of integers and booleans.
The expression \(t \mapletB u\) denotes an ordered pair.
The Cartesian product \(S \prodB T\) is a set of pairs; it should not be confused with the type product \(\tau \prodB \sigma\), although their roles are aligned.

Several familiar B connectives are treated as derived forms in BEer.
For example, disjunction, implication, equivalence, inequality, existential quantification, and binary relations can be defined from the primitive constructs:
\[
\begin{array}{rcl}
P \lor Q &\triangleq& \notB(\notB P \andB \notB Q),\\
P \Rightarrow Q &\triangleq& \notB P \lor Q,\\
P \Leftrightarrow Q &\triangleq& (P \Rightarrow Q) \andB (Q \Rightarrow P),\\
t \neq u &\triangleq& \notB(t \eqB u),\\
\exists \vec v \inB D.\,P &\triangleq& \notB(\forall \vec v \inB D.\,\notB P),\\
S \leftrightarrow T &\triangleq& (S \prodB T) \pfunB \mathbb{B}.
\end{array}
\]
Keeping these forms derived reduces the number of primitive cases required in the encoding proof.

\subsection{Typing discipline}

Typing judgments have the form \(\btype{\Gamma}{t}{\tau}\), where \(\Gamma\) maps variables to B types.
The rules below give representative cases.
Variables are typed by lookup in \(\Gamma\), and literals have their expected types:
\begin{mathpar}
\infer*[right=var]{\Gamma(x)=\tau}{\btype{\Gamma}{x}{\tau}}
\and
\infer*[right=int]{ }{\btype{\Gamma}{n}{\intB}}
\and
\infer*[right=bool]{ }{\btype{\Gamma}{b}{\boolB}}
\end{mathpar}

Pairs and arithmetic are typed structurally:
\begin{mathpar}
\infer*[right=maplet]
  {\btype{\Gamma}{t}{\tau} \quad \btype{\Gamma}{u}{\sigma}}
  {\btype{\Gamma}{t \mapletB u}{\tau \prodB \sigma}}
\and
\infer*[right=add]
  {\btype{\Gamma}{t}{\intB} \quad \btype{\Gamma}{u}{\intB}}
  {\btype{\Gamma}{t+u}{\intB}}
\and
\infer*[right=le]
  {\btype{\Gamma}{t}{\intB} \quad \btype{\Gamma}{u}{\intB}}
  {\btype{\Gamma}{t \leB u}{\boolB}}
\end{mathpar}

The boolean and equality rules are similarly standard:
\begin{mathpar}
\infer*[right=and]
  {\btype{\Gamma}{P}{\boolB} \quad \btype{\Gamma}{Q}{\boolB}}
  {\btype{\Gamma}{P \andB Q}{\boolB}}
\and
\infer*[right=not]
  {\btype{\Gamma}{P}{\boolB}}
  {\btype{\Gamma}{\notB P}{\boolB}}
\and
\infer*[right=eq]
  {\btype{\Gamma}{t}{\tau} \quad \btype{\Gamma}{u}{\tau}}
  {\btype{\Gamma}{t \eqB u}{\boolB}}
\end{mathpar}

The set rules are the ones most relevant to the later encoding:
\begin{mathpar}
\infer*[right=int-set]{ }{\btype{\Gamma}{\mathbb{Z}}{\setB \intB}}
\and
\infer*[right=bool-set]{ }{\btype{\Gamma}{\mathbb{B}}{\setB \boolB}}
\and
\infer*[right=mem]
  {\btype{\Gamma}{t}{\tau} \quad \btype{\Gamma}{S}{\setB \tau}}
  {\btype{\Gamma}{t \inB S}{\boolB}}
\and
\infer*[right=pow]
  {\btype{\Gamma}{S}{\setB \tau}}
  {\btype{\Gamma}{\powB S}{\setB(\setB \tau)}}
\and
\infer*[right=prod]
  {\btype{\Gamma}{S}{\setB \tau} \quad \btype{\Gamma}{T}{\setB \sigma}}
  {\btype{\Gamma}{S \prodB T}{\setB(\tau \prodB \sigma)}}
\and
\infer*[right=union]
  {\btype{\Gamma}{S}{\setB \tau} \quad \btype{\Gamma}{T}{\setB \tau}}
  {\btype{\Gamma}{S \cupB T}{\setB \tau}}
\and
\infer*[right=inter]
  {\btype{\Gamma}{S}{\setB \tau} \quad \btype{\Gamma}{T}{\setB \tau}}
  {\btype{\Gamma}{S \capB T}{\setB \tau}}
\end{mathpar}

Function spaces and applications are expressed through sets of pairs:
\begin{mathpar}
\infer*[right=pfun]
  {\btype{\Gamma}{S}{\setB \tau} \quad \btype{\Gamma}{T}{\setB \sigma}}
  {\btype{\Gamma}{S \pfunB T}{\setB(\setB(\tau \prodB \sigma))}}
\and
\infer*[right=app]
  {\btype{\Gamma}{f}{\setB(\tau \prodB \sigma)} \quad \btype{\Gamma}{t}{\tau}}
  {\btype{\Gamma}{f(t)}{\sigma}}
\end{mathpar}
The application rule assigns a result type whenever \(f\) is typed as a set of pairs.
Definedness is semantic: application succeeds only when that set of pairs is actually a partial function and when the argument belongs to its domain.

Binders are typed by extending the context with the bound variables.
In the formalization, a binder binds a nonempty list \(\vec v = v_1,\ldots,v_k\) whose declared types are \(\vec \tau = \tau_1,\ldots,\tau_k\).
The domain \(D\) is the Cartesian product \(D_1 \prodB \cdots \prodB D_k\), where each \(D_i\) has type \(\setB \tau_i\), and the bound variables are required to be distinct and fresh for the surrounding context.
Writing \(\Gamma[\vec v : \vec \tau]\) for the extended context and \(\vec \tau^\times\) for \(\tau_1 \prodB \cdots \prodB \tau_k\), the main binder rules are:
\begin{mathpar}
\infer*[right=all]
  {\btype{\Gamma}{D_i}{\setB\tau_i}\; (1 \le i \le k)
   \quad \btype{\Gamma[\vec v : \vec\tau]}{P}{\boolB}}
  {\btype{\Gamma}{\forall \vec v \inB D.\,P}{\boolB}}
\and
\infer*[right=collect]
  {\btype{\Gamma}{D_i}{\setB\tau_i}\; (1 \le i \le k)
   \quad \btype{\Gamma[\vec v : \vec\tau]}{P}{\boolB}}
  {\btype{\Gamma}{\{\vec v \inB D \mid P\}}{\setB\vec\tau^\times}}
\and
\infer*[right=lambda]
  {\btype{\Gamma}{D_i}{\setB\tau_i}\; (1 \le i \le k)
   \quad \btype{\Gamma[\vec v : \vec\tau]}{e}{\sigma}}
  {\btype{\Gamma}{\lambda \vec v \inB D.\,e}{\setB(\vec\tau^\times \prodB \sigma)}}
\end{mathpar}
The lambda rule again returns a set of pairs: a B lambda expression denotes a graph, not a meta-level function.
